\documentclass[11pt]{article}

\usepackage[a4paper,margin=1in]{geometry}
\usepackage{amsmath}
\usepackage{amssymb}
\usepackage{hyperref}
\usepackage{listings}

\newcommand{\dd}{\mathrm{d}}
\newcommand{\as}{\alpha_s}
\newcommand{\muf}{\mu_F}
\newcommand{\Q}{Q}
\newcommand{\gev}{\mathrm{GeV}}
\newcommand{\xf}{x f}

\title{The \texttt{proton\_pdf} Numerical Experiment}
\author{Computational HEP}
\date{}

\begin{document}
\maketitle

\section{Purpose of the experiment}

This numerical experiment is a guided tour of proton parton distribution
functions, usually called PDFs.  It does three things:
\begin{enumerate}
  \item it uses LHAPDF to plot \(x f_i(x,\muf^2)\) for several parton flavours
        and several factorisation scales \(\muf\);
  \item it checks basic proton sum rules, such as momentum conservation and
        valence-quark number conservation;
  \item it implements a deliberately small leading-order DGLAP evolution solver
        and compares its result to LHAPDF.
\end{enumerate}
The most important pedagogical goal is not to produce the most accurate PDF
evolution possible.  The goal is to make the logic visible.  We want to see
what a PDF is, why it depends on a scale, what equation controls that scale
dependence, and how one can solve that equation numerically.

\section{What is a PDF?}

A proton is not an elementary particle.  It contains quarks, antiquarks, and
gluons.  In a high-energy collision, the hard scattering usually happens over
a very short time and distance scale.  During that hard scattering, a single
parton from each proton participates directly.  The rest of the proton becomes
beam remnant.

The PDF \(f_i(x,\muf)\) is a probability-density-like object.  It tells us how
many partons of type \(i\) we expect to find inside the proton with momentum
fraction \(x\), when the proton is probed at factorisation scale \(\muf\).
The index \(i\) can be
\[
  g,\quad u,\quad \bar u,\quad d,\quad \bar d,\quad s,\quad \bar s,\quad \ldots
\]
The variable \(x\) is a fraction, so
\[
  0 < x < 1.
\]
If a proton has momentum \(P^\mu\), then a parton carrying fraction \(x\) has
approximately momentum
\[
  p^\mu \simeq x P^\mu .
\]

Strictly speaking, PDFs are not ordinary probabilities.  They are
renormalised quantum field theory matrix elements.  However, for the purposes
of a first collider-physics calculation, it is very useful to think of them as
number densities in \(x\).

\section{Why does LHAPDF return \(x f_i\)?}

LHAPDF's basic function is
\[
  \texttt{xfxQ(i,x,Q)}.
\]
It returns
\[
  x f_i(x,Q),
\]
not \(f_i(x,Q)\).  This is common in PDF plotting because \(x f_i\) has a very
direct interpretation: it is the momentum-density contribution from parton
flavour \(i\).  If \(f_i(x,Q)\dd x\) counts partons in the interval
\([x,x+\dd x]\), then
\[
  x f_i(x,Q)\dd x
\]
counts the fraction of proton momentum carried by those partons in that
interval.

That is why the experiment plots \(x f_i\).  It is also why the momentum sum
rule takes such a simple form.

\section{The factorisation scale \(\muf\)}

The symbol \(\muf\) is called the factorisation scale.  It separates two kinds
of physics:
\begin{itemize}
  \item long-distance physics inside the proton, which is put into the PDFs;
  \item short-distance physics in the hard scattering, which is computed with
        Feynman diagrams.
\end{itemize}
This separation is called factorisation.

The split is artificial.  Nature does not come with a label saying which gluon
emission is ``inside the PDF'' and which gluon emission is ``inside the hard
matrix element''.  We choose a scale \(\muf\), and we define emissions below
that scale to be absorbed into the PDF.  Emissions above that scale are treated
as part of the hard scattering.

Because the split is artificial, a fully physical prediction should not depend
on \(\muf\).  In an all-orders calculation this cancellation is exact:
\[
  \frac{\dd}{\dd\ln\muf^2}
  \left[
    \text{PDFs at }\muf
    \otimes
    \text{hard scattering at }\muf
    \right]
  =0.
\]
At finite perturbative order, the cancellation is incomplete.  That is why
varying \(\muf\) is often used to estimate missing higher-order corrections.

\section{Where does PDF running come from?}

PDFs depend on \(\muf\).  This scale dependence is often called PDF evolution
or PDF running.  The word ``running'' is used by analogy with the running of
couplings, such as \(\as(Q)\).

The physical origin is collinear radiation.  Consider a quark moving very fast
along the proton direction.  It can radiate a gluon at a very small angle:
\[
  q \to qg.
\]
Because the emission is almost parallel to the original quark, it is hard to
resolve as a separate parton unless the probe has enough resolution.  A larger
\(\muf\) means a sharper probe.  As \(\muf\) increases, the probe resolves more
of the proton's partonic substructure.  A quark at low resolution may be seen,
at higher resolution, as a quark plus a gluon.  Similarly, a gluon may split:
\[
  g \to q\bar q, \qquad g \to gg.
\]
Therefore the number and momentum distributions of quarks and gluons change
with the probing scale.

In perturbative QCD this change is calculable when \(\muf\) is large enough
that \(\as(\muf)\) is small.  The equation that describes the change is the
DGLAP equation.

\section{The DGLAP equation}

The DGLAP equation is an evolution equation in the variable
\[
  t = \ln Q^2.
\]
At leading order it reads
\[
  \frac{\partial f_i(x,Q)}{\partial \ln Q^2}
  =
  \frac{\as(Q)}{2\pi}
  \sum_j
  \int_x^1
  \frac{\dd z}{z}\,
  P_{ij}(z)\,
  f_j\left(\frac{x}{z},Q\right).
\]
Let us go through this slowly.

\begin{itemize}
  \item \(f_i(x,Q)\) is the PDF for the parton we observe after evolution.
  \item \(f_j(x/z,Q)\) is the PDF for the parent parton before the splitting.
  \item \(z\) is the momentum fraction retained by the observed daughter
        parton relative to its parent.
  \item \(P_{ij}(z)\) is a splitting kernel.  It gives the probability pattern
        for a parent parton \(j\) to produce a daughter parton \(i\).
  \item The integral from \(x\) to \(1\) says that the parent must have had at
        least as much momentum as the daughter.
\end{itemize}

The experiment evolves
\[
  F_i(x,Q) \equiv x f_i(x,Q),
\]
because this is what LHAPDF returns and what the plots show.  Multiplying the
DGLAP equation by \(x\), and using
\[
  F_j\left(\frac{x}{z},Q\right)
  =
  \frac{x}{z}
  f_j\left(\frac{x}{z},Q\right),
\]
one finds
\[
  \frac{\partial F_i(x,Q)}{\partial \ln Q^2}
  =
  \frac{\as(Q)}{2\pi}
  \sum_j
  \int_x^1
  \dd z\,
  P_{ij}(z)\,
  F_j\left(\frac{x}{z},Q\right).
\]
Notice that the explicit \(1/z\) has disappeared in the equation for \(F=x f\).
This is exactly the form implemented in the code.

\section{The leading-order splitting kernels}

The experiment uses the leading-order QCD splitting kernels.  The colour
factors for \(SU(3)\) are
\[
  C_F=\frac{4}{3}, \qquad C_A=3, \qquad T_R=\frac{1}{2}.
\]
The number of active quark flavours in the experiment is
\[
  n_f=4,
\]
corresponding to \(d,u,s,c\).

The kernels are
\[
  P_{qq}(z)
  =
  C_F
  \left[
    \frac{1+z^2}{1-z}
    \right]_+
  +
  \frac{3}{2}C_F\,\delta(1-z),
\]
\[
  P_{qg}(z)
  =
  T_R\left[z^2+(1-z)^2\right],
\]
\[
  P_{gq}(z)
  =
  C_F
  \frac{1+(1-z)^2}{z},
\]
and
\[
  P_{gg}(z)
  =
  2C_A
  \left[
    \frac{z}{(1-z)_+}
    +
    \frac{1-z}{z}
    +
    z(1-z)
    \right]
  +
  \frac{\beta_0}{2}\delta(1-z),
\]
with
\[
  \beta_0 = \frac{11C_A-4T_R n_f}{3}.
\]

The names mean the following:
\begin{itemize}
  \item \(P_{qq}\): a quark remains a quark after radiating a gluon;
  \item \(P_{qg}\): a gluon produces a quark or antiquark through
        \(g\to q\bar q\);
  \item \(P_{gq}\): a quark radiates a gluon, so the observed daughter is the
        gluon;
  \item \(P_{gg}\): a gluon remains a gluon after gluon radiation.
\end{itemize}

\section{What is the plus prescription?}

Some splitting kernels contain expressions like
\[
  \frac{1}{1-z}.
\]
This becomes singular when \(z\to1\).  Physically, \(z\to1\) corresponds to an
emission in which the emitted parton is very soft.  The real-emission
singularity is cancelled by virtual corrections.  In the compact DGLAP kernels,
this cancellation is encoded by the plus prescription.

For a singular function \(g(z)\), the plus distribution is defined by how it
acts under an integral.  If \(h(z)\) is a smooth test function, then
\[
  \int_0^1 \dd z\, [g(z)]_+ h(z)
  =
  \int_0^1 \dd z\, g(z)\,[h(z)-h(1)].
\]
The subtraction \(h(z)-h(1)\) makes the integral finite at \(z=1\).

Our convolution does not start at \(0\).  It starts at \(x\):
\[
  \int_x^1 \dd z\, [g(z)]_+ h(z).
\]
The code uses the identity
\[
  \int_x^1 \dd z\, [g(z)]_+ h(z)
  =
  \int_x^1 \dd z\, g(z)\,[h(z)-h(1)]
  -
  h(1)\int_0^x \dd z\, g(z).
\]
For the \(P_{qq}\) singular piece,
\[
  g_{qq}(z)=C_F\frac{1+z^2}{1-z},
\]
the primitive needed for the second term is
\[
  \int_0^x \dd z\, g_{qq}(z)
  =
  C_F
  \left[
    -2\ln(1-x)-x-\frac{x^2}{2}
    \right].
\]
For the plus part of \(P_{gg}\),
\[
  g_{gg}(z)=2C_A\frac{z}{1-z},
\]
the primitive is
\[
  \int_0^x \dd z\, g_{gg}(z)
  =
  2C_A[-\ln(1-x)-x].
\]
These formulae are implemented explicitly.  This is a very useful part of the
exercise because most production PDF-evolution codes hide the plus-prescription
bookkeeping from the user.

\section{The evolution system solved by the code}

The state vector contains
\[
  F_g(x),\quad
  F_q(x),\quad
  F_{\bar q}(x)
\]
for the active flavours \(q=d,u,s,c\).  At a fixed \(x\), the quark equation is
schematically
\[
  \frac{\partial F_q}{\partial\ln Q^2}
  =
  \frac{\as}{2\pi}
  \left[
    P_{qq}\otimes F_q
    +
    P_{qg}\otimes F_g
    \right],
\]
and the antiquark equation is the same:
\[
  \frac{\partial F_{\bar q}}{\partial\ln Q^2}
  =
  \frac{\as}{2\pi}
  \left[
    P_{qq}\otimes F_{\bar q}
    +
    P_{qg}\otimes F_g
    \right].
\]
The gluon equation is
\[
  \frac{\partial F_g}{\partial\ln Q^2}
  =
  \frac{\as}{2\pi}
  \left[
    P_{gg}\otimes F_g
    +
    \sum_q P_{gq}\otimes(F_q+F_{\bar q})
    \right].
\]
Here \(\otimes\) denotes the convolution
\[
  (P\otimes F)(x)
  =
  \int_x^1 \dd z\,P(z)\,F\left(\frac{x}{z}\right).
\]

\section{Discretising the \(x\) dependence}

The DGLAP equation is not an ordinary differential equation in one variable.
It is an integro-differential equation.  The derivative is with respect to
\(\ln Q^2\), but the right-hand side contains an integral over \(z\) and the
PDF evaluated at shifted points \(x/z\).

The experiment turns it into a finite-dimensional ODE problem as follows.

\begin{enumerate}
  \item Choose an \(x\)-grid:
        \[
          x_1,x_2,\ldots,x_N.
        \]
        The code uses a hybrid grid, with many points at small \(x\) and a
        linear tail toward large \(x\).  This is because PDFs often change very
        quickly at small \(x\), while valence structure is important at
        moderate \(x\).
  \item Store the values
        \[
          F_i(x_a,Q)
        \]
        for every parton flavour \(i\) and every grid point \(x_a\).
  \item When the convolution asks for
        \[
          F_i(x_a/z,Q),
        \]
        the point \(x_a/z\) will usually not lie exactly on the grid.  The code
        therefore uses linear interpolation.
  \item At \(x=1\), the PDF is forced to zero.  This is physically sensible:
        a single parton carrying exactly all of the proton momentum has zero
        density in the continuous distribution.
  \item The \(z\)-integrals are evaluated by a trapezoidal quadrature.
\end{enumerate}

After these steps, the DGLAP equation becomes a large but ordinary system:
\[
  \frac{\dd \mathbf{F}}{\dd t}
  =
  \mathbf{R}(\mathbf{F},t),
  \qquad
  t=\ln Q^2.
\]
The vector \(\mathbf{F}\) contains all parton flavours at all \(x\)-grid points.
The function \(\mathbf{R}\) is the right-hand side built from the splitting
kernels and convolutions.

\section{The Runge--Kutta scale evolution}

To evolve from \(Q_0\) to \(Q_1\), the code evolves in
\[
  t=\ln Q^2.
\]
The starting and ending values are
\[
  t_0 = \ln Q_0^2,
  \qquad
  t_1 = \ln Q_1^2.
\]
If the user requests \(N_{\mathrm{step}}\) evolution steps, then
\[
  \Delta t = \frac{t_1-t_0}{N_{\mathrm{step}}}.
\]

The code uses the fourth-order Runge--Kutta method.  For an ODE
\[
  \frac{\dd \mathbf{F}}{\dd t} = \mathbf{R}(\mathbf{F},t),
\]
one step is
\[
  \mathbf{k}_1 = \mathbf{R}(\mathbf{F}_n,t_n),
\]
\[
  \mathbf{k}_2 =
  \mathbf{R}\left(
  \mathbf{F}_n+\frac{\Delta t}{2}\mathbf{k}_1,
  t_n+\frac{\Delta t}{2}
  \right),
\]
\[
  \mathbf{k}_3 =
  \mathbf{R}\left(
  \mathbf{F}_n+\frac{\Delta t}{2}\mathbf{k}_2,
  t_n+\frac{\Delta t}{2}
  \right),
\]
\[
  \mathbf{k}_4 =
  \mathbf{R}\left(
  \mathbf{F}_n+\Delta t\,\mathbf{k}_3,
  t_n+\Delta t
  \right),
\]
and
\[
  \mathbf{F}_{n+1}
  =
  \mathbf{F}_n
  +
  \frac{\Delta t}{6}
  \left(
  \mathbf{k}_1
  +2\mathbf{k}_2
  +2\mathbf{k}_3
  +\mathbf{k}_4
  \right).
\]
This is a standard ODE method.  It is not specific to QCD.  The QCD content is
inside the right-hand side \(\mathbf{R}\).

\section{Why the RK4 formula is reasonable}

It is useful to justify the Runge--Kutta formula, because otherwise it can look
like a mysterious recipe.  The starting point is the exact integral form of the
ODE.  If
\[
  \frac{\dd \mathbf{F}}{\dd t} = \mathbf{R}(\mathbf{F},t),
\]
then over one small step,
\[
  \mathbf{F}(t_n+\Delta t)
  =
  \mathbf{F}(t_n)
  +
  \int_{t_n}^{t_n+\Delta t}
  \dd t\,
  \mathbf{R}(\mathbf{F}(t),t).
\]
So the whole problem is this: we need the average value of the slope
\(\mathbf{R}\) across the interval, but we do not know the exact curve
\(\mathbf{F}(t)\) inside the interval.

The simplest method, Euler's method, uses only the slope at the beginning:
\[
  \mathbf{F}_{n+1}
  \approx
  \mathbf{F}_n
  +
  \Delta t\,\mathbf{R}(\mathbf{F}_n,t_n).
\]
This is easy to understand, but it is crude.  It assumes that the slope at the
left endpoint is representative of the whole interval.  If the slope changes
noticeably during the step, Euler's method misses that change.

RK4 improves this by sampling the slope four times:
\begin{itemize}
  \item \(\mathbf{k}_1\) is the slope at the beginning of the step;
  \item \(\mathbf{k}_2\) is an estimate of the slope halfway through the step,
        using \(\mathbf{k}_1\) to guess the midpoint value of \(\mathbf{F}\);
  \item \(\mathbf{k}_3\) is another midpoint slope, now using the better
        midpoint estimate from \(\mathbf{k}_2\);
  \item \(\mathbf{k}_4\) is an estimate of the slope at the end of the step,
        using \(\mathbf{k}_3\) to guess the endpoint value of \(\mathbf{F}\).
\end{itemize}
The final combination
\[
  \frac{1}{6}
  \left(
  \mathbf{k}_1
  +2\mathbf{k}_2
  +2\mathbf{k}_3
  +\mathbf{k}_4
  \right)
\]
is a weighted average slope.  The midpoint slopes receive twice the weight of
the endpoint slopes because they usually give a better estimate of the average
slope across the interval.  This is closely related in spirit to Simpson's rule
for numerical integration, where midpoint information greatly improves the
estimate of an integral.

One can make this statement precise by expanding the exact solution in a Taylor
series in \(\Delta t\).  RK4 is constructed so that its Taylor expansion agrees
with the exact solution through terms of order \((\Delta t)^4\).  Therefore the
error made in a single step is of order
\[
  (\Delta t)^5,
\]
and after many steps over a fixed interval the accumulated error is of order
\[
  (\Delta t)^4.
\]
That is why it is called a fourth-order method.  For this exercise, RK4 is a
good compromise: it is much more accurate than Euler's method, but it is still
short enough to write explicitly and inspect line by line.

At every intermediate RK4 evaluation, the code converts the current \(t\) back
to a scale
\[
  Q = e^{t/2}
\]
and asks LHAPDF for \(\as(Q)\).  Thus the pedagogical DGLAP solver uses
LHAPDF's running coupling but computes the PDF evolution itself.

The plot produced by the experiment now shows two RK4 choices.  The first is
the user-configured evolution in equal steps of \(t=\ln Q^2\).  The second is a
fixed comparison evolution in ten equal steps of \(Q\) itself.  For the default
comparison
\[
  Q_0=2~\gev,\qquad Q_1=10~\gev,
\]
the second path uses
\[
  Q = 2.0,\ 2.8,\ 3.6,\ \ldots,\ 10.0~\gev.
\]
Each of these intervals is still evolved with RK4 in the variable
\(t=\ln Q^2\); only the way the interval is subdivided is different.  Showing
both curves is a useful numerical check: if the step size is small enough and
the implementation is stable, changing the subdivision should not drastically
change the result.
The bottom panel of the comparison plot makes this explicit by showing
\[
  \log_{10}
  \left(
    \frac{|F_{\mathrm{uniform}\ \ln Q^2}(x)-F_{\mathrm{linear}\ Q}(x)|}
         {|F_{\mathrm{LHAPDF}}(x,Q_1)|}
  \right),
\]
so that two visually overlapping curves can still be checked quantitatively.

\section{Where Symbolica enters}

The code uses Symbolica to build readable symbolic expressions for the
splitting kernels.  For example, the singular part of \(P_{qq}\) is represented
symbolically as
\[
  C_F\frac{1+z^2}{1-z}.
\]
The numerical quadrature then uses fast vectorised functions corresponding to
those symbolic expressions.  This gives a useful compromise: the code contains
explicit mathematical formulae for the kernels, but the default run is still
fast enough for a classroom exercise.

\section{Comparison with LHAPDF}

The DGLAP comparison proceeds as follows:
\begin{enumerate}
  \item choose a starting scale \(Q_0\);
  \item query LHAPDF for \(F_i(x,Q_0)=x f_i(x,Q_0)\) on the \(x\)-grid;
  \item evolve those values to a target scale \(Q_1\) with the pedagogical LO
        solver, once with the configured equal-\(\ln Q^2\) RK4 steps and once
        with ten equal-\(Q\) RK4 steps;
  \item query LHAPDF directly at \(Q_1\);
  \item plot the starting LHAPDF shape, the target LHAPDF shape, both evolved
        answers, and the relative differences.
\end{enumerate}
The comparison plot is named
\[
  \texttt{proton\_pdf\_dglap\_lo\_vs\_lhapdf.png}.
\]

The evolved curves should not be expected to agree perfectly with LHAPDF.  The
comparison is not a precision validation of a full PDF-evolution program.
Differences arise because:
\begin{itemize}
  \item the experiment uses leading-order splitting kernels;
  \item the default LHAPDF set is an NLO fitted PDF set;
  \item the \(x\)-grid is coarse compared with production codes;
  \item interpolation and trapezoidal quadrature are simple approximations;
  \item heavy-flavour threshold treatment is simplified;
  \item only the active flavours \(d,u,s,c\) are evolved in the toy solver.
\end{itemize}
The point is to see the mechanism of PDF running, not to reproduce LHAPDF's
internal machinery exactly.

\section{Momentum and valence sum rules}

The experiment also checks several sum rules.  These are conservation laws or
proton quantum-number statements written in PDF language.

\subsection{Momentum sum rule}

All partons together carry all of the proton momentum.  Therefore
\[
  \int_0^1 \dd x\,
  \left[
    xg(x,Q)
    +
    \sum_q xq(x,Q)
    +
    \sum_q x\bar q(x,Q)
    \right]
  =1.
\]
Because LHAPDF returns \(x f_i\), the integrand is directly the sum of LHAPDF
queries:
\[
  xg + xq + x\bar q.
\]
The code performs this integral numerically on the \(x\)-grid.

\subsection{Valence sum rules}

The proton contains two valence up quarks and one valence down quark.  Sea
quarks come in quark-antiquark pairs and therefore do not contribute to valence
number.  The valence distributions are
\[
  u_v(x,Q)=u(x,Q)-\bar u(x,Q),
\]
\[
  d_v(x,Q)=d(x,Q)-\bar d(x,Q).
\]
The sum rules are
\[
  \int_0^1 \dd x\,u_v(x,Q)=2,
  \qquad
  \int_0^1 \dd x\,d_v(x,Q)=1.
\]
For strange and charm in this simplified check,
\[
  \int_0^1 \dd x\,s_v(x,Q)=0,
  \qquad
  \int_0^1 \dd x\,c_v(x,Q)=0.
\]
Here the code must divide LHAPDF's \(xq\) by \(x\), because the valence number
integral is over \(q-\bar q\), not over \(x(q-\bar q)\):
\[
  q_v(x,Q)
  =
  \frac{xq(x,Q)-x\bar q(x,Q)}{x}.
\]

\subsection{Baryon number}

Each quark carries baryon number \(1/3\).  Antiquarks carry baryon number
\(-1/3\).  The proton has baryon number 1.  Therefore
\[
  B
  =
  \frac{1}{3}
  \int_0^1 \dd x\,
  \left[
    u_v(x,Q)+d_v(x,Q)+s_v(x,Q)+c_v(x,Q)
    \right]
  \simeq 1.
\]

\subsection{Isospin}

For the proton, the third component of isospin is
\[
  I_3 = +\frac{1}{2}.
\]
In this simplified parton-counting picture,
\[
  I_3
  =
  \frac{1}{2}
  \left[
    \int_0^1 \dd x\,u_v(x,Q)
    -
    \int_0^1 \dd x\,d_v(x,Q)
    \right]
  \simeq \frac{1}{2}.
\]

\section{Numerical integration of the sum rules}

The sum rules are integrals over \(x\in[0,1]\).  The code cannot evaluate a
continuum integral exactly.  It samples a finite \(x\)-grid and uses the
trapezoidal rule:
\[
  \int_a^b \dd x\,F(x)
  \approx
  \sum_{k=1}^{N-1}
  \frac{x_{k+1}-x_k}{2}
  \left[
    F(x_k)+F(x_{k+1})
    \right].
\]
The resulting numbers are written to
\[
  \texttt{proton\_pdf\_sum\_rules.csv}.
\]
Small deviations from the exact targets can come from grid truncation, grid
coarseness, and interpolation details in the PDF set.

\section{Example command}

A typical run is
\begin{lstlisting}[language=bash]
./run.py --proton_pdf \
  --pdf_set NNPDF40_nlo_as_01180_nf_4 \
  --q_values 2.0,10.0,91.1876 \
  --flavours g,u,d,s,c,ubar,dbar,sbar,cbar \
  --dglap_start_q 2.0 \
  --dglap_target_q 10.0 \
  --dglap_flavour u \
  --dglap_x_points 45 \
  --dglap_z_points 80 \
  --dglap_steps 4 \
  --output_dir proton_pdf_lo_run
\end{lstlisting}
The output directory contains:
\[
  \texttt{proton\_pdf\_xfx\_Q\_*.png},
\]
\[
  \texttt{proton\_pdf\_sum\_rules.csv},
\]
\[
  \texttt{proton\_pdf\_dglap\_lo\_vs\_lhapdf.png}.
\]

\section{What to learn from this experiment}

The main lessons are:
\begin{itemize}
  \item PDFs describe the proton's partonic content as a function of momentum
        fraction \(x\) and scale \(Q\).
  \item The scale dependence is not arbitrary.  It is caused by resolvable
        collinear QCD radiation.
  \item DGLAP evolution turns this physical idea into an equation.
  \item The equation is non-local in \(x\): the PDF at \(x\) receives
        contributions from parent partons at \(x/z\).
  \item The plus prescription is the mathematical way real and virtual soft
        singularities combine into a finite evolution equation.
  \item Numerical PDF evolution requires choices: an \(x\)-grid, interpolation,
        quadrature, a running coupling, and an ODE solver.
  \item Sum rules are powerful checks that the PDFs still describe a proton.
\end{itemize}

\end{document}
